\documentclass{article}

\usepackage[utf8]{inputenc}
\usepackage[T1]{fontenc}
\usepackage{amsmath}
\usepackage{amsfonts}
\usepackage{amssymb}
\usepackage{graphicx}
\usepackage{fourier}
\usepackage[hidelinks]{hyperref}
\usepackage{indentfirst}

%opening
\title{NLO Equations}
\author{Simone Florincy Torres Marques}

\begin{document}
	
	\maketitle
	
	\begin{abstract}
	\end{abstract}
	
	\section{Electric Dipole}
	\begin{displaymath}
	\mu = \sqrt{{\mu_x}^2+{\mu_y}^2+{\mu_z}^2}
	\end{displaymath}
	
	\section{Polarizability}
	\begin{displaymath}
	\alpha = \frac{{\mathrm{xx}} + {\mathrm{yy}} + {\mathrm{zz}}}{3} 
	\end{displaymath}
	
	\section{First Hyperpolarizability}
		\paragraph{}
	The tensor components $\beta_i$ generated by Gaussian lack multiplication by a factor of $\frac{1}{3}$. Additionally, the conversion from the Taylor series convention to the perturbational series convention involves multiplying by a $\frac{1}{2}$ factor, when considering the 3 possible beta tensors founded in Gaussian 16.0. This results in a total factor of $\frac{1}{6}$ to obtain $\beta$ in the latter convention, as illustrated below:
	\begin{align*}
	\beta_i &=\beta_{iii}+\frac {1} {3}{\sum_{i\neq j}}(\beta_{ijj}+\beta_{jij}+\beta_{jji})\\
	\beta_i^{G16} &=3\beta_i
	\end{align*}
	\begin{align*}
	\beta^{B} &= \frac{1}{2} \beta^{T} \\
	\beta^{B} &=  \frac{1}{6}\beta^{G16}
	\end{align*}
	\subsection{T convention}
	\begin{displaymath}
	\beta = \sqrt{\left(\frac{x}{3}\right)^2 + \left(\frac{y}{3}\right)^2 + \left(\frac{z}{3}\right)^2} 
	\end{displaymath}
	\subsection{B convention}
	\begin{displaymath}
	\beta = \sqrt{\left(\frac{x}{6}\right)^2 + \left(\frac{y}{6}\right)^2 + \left(\frac{z}{6}\right)^2} 
	\end{displaymath}
	\section{Second Hyperpolarizability}
	\paragraph{}
	The tensor components generated by Gaussian lack multiplication by a factor of $\frac{1}{6}$. Additionally, the conversion from the Taylor series convention to the perturbational series convention involves multiplying by a $\frac{1}{6}$ factor, when considering the 3 possible gamma tensors founded in Gaussian 16.0. This results in a total factor of $\frac{1}{36}$ to obtain $\gamma$ in the latter convention, as illustrated below:
	\begin{align*}
	\frac{3}{2} \gamma^{B} &= \frac{1}{4} \gamma^{T} \\
	\gamma^{T} &= \frac{1}{6} \gamma^{G16} \\
	\gamma^{B} &= \frac{1}{36}\gamma^{G16} 
	\end{align*}
	\paragraph{}$\gamma$ represents the overall gamma property calculated as discussed in the following sections.
	\subsection{Gamma Static (0;0,0,0)}
	\subsubsection{T convention}
	\begin{displaymath}
	\gamma = \frac{1}{5}\left(\frac{\mathrm{xxxx}}{6} + \frac{\mathrm{yyyy}}{6} + \frac{\mathrm{zzzz}}{6} + 2\frac{\mathrm{xxyy}}{6} + 2\frac{\mathrm{xxzz}}{6} + 2\frac{\mathrm{yyzz}}{6}\right)
	\end{displaymath}
	\subsubsection{B convention}
	\begin{displaymath}
	\gamma = \frac{1}{5}\left(\frac{\mathrm{xxxx}}{36} + \frac{\mathrm{yyyy}}{36} + \frac{\mathrm{zzzz}}{36} + 2\frac{\mathrm{xxyy}}{36} + 2\frac{\mathrm{xxzz}}{36} + 2\frac{\mathrm{yyzz}}{36}\right)
	\end{displaymath}
	\subsection{Gamma Kerr effect (-w;w,0,0)}
	\subsubsection{T convention}
	\begin{displaymath}
	\gamma = \frac{1}{5}\left(\frac{{xxxx}}{6} + \frac{{yyyy}}{6} + \frac{{zzzz}}{6} + \frac{{xxyy}}{6} + \frac{{yyxx}}{6} + \frac{{xxzz}}{6} + \frac{{zzxx}}{6} + \frac{{yyzz}}{6} + \frac{{zzyy}}{6}\right)
	\end{displaymath}
	\subsubsection{B convention}
	\begin{displaymath}
	\gamma = \frac{1}{5}\left(\frac{{xxxx}}{36} + \frac{{yyyy}}{36} + \frac{{zzzz}}{36} + \frac{{xxyy}}{36} + \frac{{yyxx}}{36} + \frac{{xxzz}}{36} + \frac{{zzxx}}{36} + \frac{{yyzz}}{36} + \frac{{zzyy}}{36}\right)
	\end{displaymath}
	\subsection{Gamma EFISH (-2w,w,w,0)}
	\subsubsection{T convention}
	\begin{align*}
	\gamma &= 3\left(\frac{\mathrm{xxxx}}{6} + \frac{\mathrm{yyyy}}{6} + \frac{\mathrm{zzzz}}{6}\right) + 2\left(\frac{\mathrm{xyxy}}{6} + \frac{\mathrm{xzxz}}{6} + \frac{\mathrm{yzyz}}{6} + \frac{\mathrm{yxyx}}{6} + \frac{\mathrm{zxzx}}{6} + \frac{\mathrm{zyzy}}{6}\right) \nonumber \\
	&+ \frac{\mathrm{xyyx}}{6} + \frac{\mathrm{xzzx}}{6} + \frac{\mathrm{yxxy}}{6} + \frac{\mathrm{yzzy}}{6} + \frac{\mathrm{zxxz}}{6} + \frac{\mathrm{zyyz}}{6}  
	\end{align*}
	
\end{document}
\subsubsection{B convention}
	\begin{align*}
	\gamma &= 3\left(\frac{\mathrm{xxxx}}{36} + \frac{\mathrm{yyyy}}{36} + \frac{\mathrm{zzzz}}{36}\right) + 2\left(\frac{\mathrm{xyxy}}{36} + \frac{\mathrm{xzxz}}{36} + \frac{\mathrm{yzyz}}{36} + \frac{\mathrm{yxyx}}{36} + \frac{\mathrm{zxzx}}{36} + \frac{\mathrm{zyzy}}{36}\right) \nonumber \\
	&+ \frac{\mathrm{xyyx}}{36} + \frac{\mathrm{xzzx}}{36} + \frac{\mathrm{yxxy}}{36} + \frac{\mathrm{yzzy}}{36} + \frac{\mathrm{zxxz}}{36} + \frac{\mathrm{zyyz}}{36}  
	\end{align*}